% Generated from locale/tr/content/proof-theory/propositions-as-types/terms-to-proofs.tex; do not hand-edit.


\olfileid[tr]{int}{pty}{tp}

\olsection{İspat Terimlerinden \usetoken{P}{derivation}{}i Geri Elde Etme}

Şimdi öteki yönü, yani bir ispat terimi~$N$'yi \Log{N2i}'deki bir
!!a{proof}{}a geri çevirmeyi ele alalım. Genel olarak bu ancak, verilen ispat
teriminde serbest geçen her~$x$ değişkeni için $x : !A$ çiftini içeren bir
bağlama göre yapılabilir. Ancak serbest değişken bulunmayan şu terimle
başlayalım:
\[
\lambd[\typeof{x}{!A \land
  !B}][\pair{\proj{2}{x}}{\proj{1}{x}}]
\]
Bu terim $\lambd[\typeof{x}{!A \land !B}][N_1]$ biçimindedir. \Log{tN2i}'de
bu biçimde bir terim üreten tek kural~\Intro{\lif} olduğundan $N$'nin
kodladığı !!{proof}{}ın \Intro{\lif} ile bitmek zorunda olduğunu biliriz;
yani şöyle biter:
\[
  \Axiom$x : !A\land !B \fCenter \pair{\proj{2}{x}}{\proj{1}{x}} : 
    !C$
  \RightLabel{\Intro\lif}
  \UnaryInf$\fCenter \lambd[\typeof{x}{!A \land
  !B}][\pair{\proj{2}{x}}{\proj{1}{x}}] : (!A \land !B) \lif !C$
  \DisplayProof
\]
$!C$'nin ne olduğunu henüz bilmiyoruz; ancak önbileşenin $!A \land !B$
olduğunu ve öncül bağlamının $x:!A \land !B$'yi içermek zorunda olduğunu
biliyoruz.

Şimdi öncülle ilişkili ispat terimi zaten belirlenmiştir:
$\pair{\proj{2}{x}}{\proj{1}{x}}$. Onun kodladığı !!{proof},
\Intro{\land} ile bitmek zorundadır. $!C$'nin ne olduğunu hâlâ bilmiyoruz;
ancak \Intro{\land}'ın sonucu olduğu için ve bağlaçlı olması,
yani $!C \ident (!C_1 \land !C_2)$ olması gerektiğini biliyoruz. Yeniden
kurmayı sürdürürüz:
\[
  \Axiom$x : !A\land !B \fCenter \proj{2}{x} : !C_1$
  \Axiom$x : !A\land !B \fCenter \proj{1}{x} : !C_2$
  \RightLabel{\Intro\land}
  \BinaryInf$x : !A\land !B \fCenter \pair{\proj{2}{x}}{\proj{1}{x}} : 
    !C_1 \land !C_2$
  \RightLabel{\Intro\lif}
  \UnaryInf$\fCenter \lambd[\typeof{x}{!A \land
  !B}][\pair{\proj{2}{x}}{\proj{1}{x}}] : (!A \land !B) \lif (!C_1 \land !C_2)$
  \DisplayProof
\]
İspat terimi $\proj{2}{x}$, soldaki üst sekantın \Elim{\land}[2] ile elde
edildiği, yani öncülün $!D \land !C_1$ biçiminde olduğu anlamına gelir.
Sağda $!C_2$'ye götüren çıkarımın \Elim{\land}[1] olması ve öncülün
$!C_2 \land !E$ biçiminde bulunması gerektiğini belirleyebiliriz:
\[
  \Axiom$x : !A\land !B \fCenter x: !D\land !C_1$
  \RightLabel{\Elim\land[2]}
  \UnaryInf$x : !A\land !B \fCenter \proj{2}{x} : !C_1$
  \Axiom$x : !A\land !B \fCenter x: !C_2\land !E$
  \RightLabel{\Elim\land[1]}
  \UnaryInf$x : !A\land !B \fCenter \proj{1}{x} : !C_2$
  \RightLabel{\Intro\land}
  \BinaryInf$x : !A\land !B \fCenter \pair{\proj{2}{x}}{\proj{1}{x}} : 
    !C_1 \land !C_2$
  \RightLabel{\Intro\lif}
  \UnaryInf$\fCenter \lambd[\typeof{x}{!A \land
  !B}][\pair{\proj{2}{x}}{\proj{1}{x}}] : (!A \land !B) \lif (!C_1 \land !C_2)$
  \DisplayProof
\]
En üstteki sekantların ispat terimleri artık değişken olduğundan bunların
aksiyom olması gerektiğini biliriz. Bu, $!C_1$, $!C_2$, $!D$ ve $!E$'yi
belirler: $!D \ident !A$ ve $!C_1 \ident !B$ olmalıdır; yoksa sol sekant
aksiyom değildir. Benzer biçimde $!C_2 \ident !A$ ve $!E \ident !B$
olmalıdır; yoksa sağ sekant aksiyom değildir. İspatı bütünüyle geri elde
etmiş olduk:
\[
  \Axiom$x : !A\land !B \fCenter x: !A\land !B$
  \RightLabel{\Elim\land[2]}
  \UnaryInf$x : !A\land !B \fCenter \proj{2}{x} : !B$
  \Axiom$x : !A\land !B \fCenter x: !A\land !B$
  \RightLabel{\Elim\land[1]}
  \UnaryInf$x : !A\land !B \fCenter \proj{1}{x} : !A$
  \RightLabel{\Intro\land}
  \BinaryInf$x : !A\land !B \fCenter \pair{\proj{2}{x}}{\proj{1}{x}} : 
    !B \land !A$
  \RightLabel{\Intro\lif}
  \UnaryInf$\fCenter \lambd[\typeof{x}{!A \land
  !B}][\pair{\proj{2}{x}}{\proj{1}{x}}] : (!A \land !B) \lif (!B \land !A)$
  \DisplayProof
\]


